\section{QA Skills ובקרת איכות}

ייצור \eng{PDF} אוטומטי במסמכים ארוכים (30--50 עמודים) מייצר באגים חוזרים:
טבלאות שבורות, עברית-אנגלית מעורבבת, תמונות חורגות מהשוליים. המרצה
הציג מערכת \eng{QA Skills} היררכית לטיפול בכך --- מודל שניתן ליישם
בכל פרויקט סוכן.

% QA orchestration hierarchy
\begin{figure}[H]
\centering
\begin{tikzpicture}[
  level 1/.style={sibling distance=4.2cm, level distance=1.6cm},
  level 2/.style={sibling distance=2.2cm, level distance=1.4cm},
  every node/.style={draw, rounded corners=2pt, align=center, font=\footnotesize,
    minimum width=2cm, minimum height=0.7cm}
]
  \node[fill=primary!15, draw=primary, font=\bfseries] {QA Super}
    child { node[fill=accentblue!12] {QA Table} }
    child { node[fill=accentgreen!12] {QA Bidi} }
    child { node[fill=accentamber!12] {QA Image} }
    child { node[fill=lightgray] {QA Code} };
\end{tikzpicture}
\caption{היררכיית \eng{QA Skills} לפי דוגמת המרצה: מנהל-על ומומחי תחום}
\label{fig:qa-hierarchy}
\end{figure}


\subsection{פילוסופיית QA}

כאשר מתגלה בעיה:
\begin{enumerate}[rightmargin=2em]
  \item \textbf{אל תתקן סתם} --- בקש הסבר מה קרה.
  \item החלט: האם הפתרון שייך ל-\eng{Skill} (שפה טבעית) או ל-\eng{CLS} (טכני)?
  \item עדכן את השכבה המתאימה.
  \item הרץ מחדש ובדוק שהתיקון כללי.
\end{enumerate}

\subsection{מבנה היררכי}

לפי דוגמת המרצה:
\begin{description}[rightmargin=2em, style=nextline]
  \item[QA Super] ``מנכ''ל'' --- מכיר את כל סוגי הבעיות, מפעיל מומחים.
  \item[QA Table] מחלקת טבלאות --- יישור, עמודות, עברית-אנגלית בטבלה.
  \item[QA Bidi] כיווניות --- RTL/LTR, מספרים, נוסחאות.
  \item[QA Image] תמונות --- גבולות, כיתובים, מסגרות.
  \item[QA Code] קטעי קוד --- רקע, גלישת שורות, התנגשויות.
\end{description}

כל ``מחלקה'' היא \eng{Skill} (או סוכן עם \eng{Skill} ייעודי) עם כללי
זיהוי (\eng{detector}) ותיקון (\eng{fixer}).

\subsection{Detector ו-Fixer}

\begin{itemize}[rightmargin=2em]
  \item \textbf{Detector} --- רץ על \eng{PDF} או על קוד מקור, מדווח בעיות
    (למשל: ``הטבלה זזה ימינה'').
  \item \textbf{Fixer} --- מקבל דיווח, יודע לתקן (למשל: עדכון הגדרות
    עמודה ב-\eng{CLS}).
\end{itemize}

\subsection{דוגמת Skill QA-Table}

\begin{lstlisting}[caption={קטע מ-QA Skill לטבלאות}]
# QA-Table Skill
Detect:
- Tables aligned to wrong margin in RTL documents
- Column order reversed (Hebrew headers)
Fix:
- Apply \begin{RTL} inside table environment
- Use p{} columns with \RaggedLeft
- Update agentic-ai-template.cls @AtBeginEnvironment{table}
Never: patch only the single table in main.tex
\end{lstlisting}

\subsection{יתרונות מערכת QA}

\agenttable{
\caption{רכיבי מערכת \eng{QA} למסמכי \eng{PDF}}
\label{tab:qa-system}
\begin{tabular}{@{}p{3cm}p{4.5cm}p{4.5cm}@{}}
\toprule
\textbf{רכיב} & \textbf{תפקיד} & \textbf{פלט} \\
\midrule
Detector & זיהוי חריגה & דוח בעיות \\
Fixer & תיקון ממוקד & \eng{diff} ב-\eng{CLS}/\eng{Skill} \\
QA Super & תיאום & אישור לפני הגשה \\
Compile & אימות & \eng{PDF} סופי \\
\bottomrule
\end{tabular}
}

\subsection{יישום במטלה 04}

רשימת בדיקות לפני הגשה:
\begin{itemize}[rightmargin=2em]
  \item קומפילציה ללא שגיאות (\texttt{compile.ps1}).
  \item תוכן עניינים, רשימת איורים וטבלאות מעודכנים.
  \item ביבליוגרפיה עם ציטוטים [1], [2] בגוף המסמך.
  \item לפחות 30 עמודים, 5 טבלאות, 5 איורים.
  \item עברית-אנגלית ללא שבירות בולטות.
\end{itemize}

\begin{insightbox}
\eng{QA Skills} הם ``בטיחות בדם נקנית'' --- כל שורה בקובץ ה-\eng{Skill}
מייצגת באג אמיתי מהעבר. עם הזמן הספרייה הופכת יקרת ערך.
\end{insightbox}
